\section{Appendix 21: Scenario Phi - Elder Care}

\subsection{The Legacy Paradigm \\\& Its Failures}
Legacy elder care technology relies on intrusive monitoring, complex mobile apps with steep learning curves, and centralized cloud servers that process sensitive health data. Notifications cause anxiety, and the interfaces are fundamentally inaccessible to those with cognitive or physical decline, treating them as data subjects rather than autonomous individuals.

\subsection{The Incremental Automation Pathway}
Phase 1: Caregivers manually log health metrics into shared local ledgers.
Phase 2: Ambient sensors and wearables continuously broadcast telemetry via the local mesh, subtly integrating with spatial UIs (e.g., smart mirrors) to gently guide daily routines.
Phase 3: The system reaches zero-click operation, automatically adjusting environmental parameters (lighting, temperature) and anticipating needs without explicit user commands.

\subsection{The Human Boundary (Limits of Automation)}
Algorithmic adjustments cannot replace human connection. The interface must pause and mandate human interaction for emotional support or significant medical interventions. The system will withhold automatic medication adjustments if the patient's emotional state dictates that a trusted human caregiver should deliver the news or the care directly.

\subsection{Interface Mechanics (The "How")}
This scenario utilizes spatial and ambient interfaces heavily, minimizing direct screen time. Data is processed locally to preserve dignity and privacy. The Cognitive Load algorithm acts as a "Centaur" interface for caregivers, aggregating alerts into a daily digest rather than a barrage of push notifications, utilizing local CRDTs to sync across the caregiving team.
